\section{The renormalized variational problem}\label{sec:variational}

Set $q=a/2$. We first define an exact action on geodesic polygons, then
relax it in the uniform topology. The construction will not assume that
a limiting norm is split.

\subsection{Exact polygonal action}
Let $I=[l,r]$ and fix ordered endpoint pairs. An ordered polygon $c$ is
a continuous curve in $\Dord$ which is an affinely parametrized product
geodesic on each interval of a finite partition
$l=u_0<\cdots<u_N=r$. Its nodes may be arbitrary norms. Write
\[
 \Delta_j=u_{j+1}-u_j,\qquad
 W_j=\int_{u_j}^{u_{j+1}}R(u)\,\dd u,\qquad
 g(c)=d(c_C)-d(c_A)\ge0.
\]
Its action is
\begin{equation}\label{eq:polygonaction}
 E_I(c)=\sum_j\left\{
 \frac{W_j}{2}\sum_{v=A,C}\sum_i
   \ell\left(\frac{\sigma_i(c_v(u_j),c_v(u_{j+1}))}{\Delta_j}\right)
 +q\int_{u_j}^{u_{j+1}}R(u)g(c(u))\,\dd u\right\}.
\end{equation}
The determinant gap is affine on a geodesic. Its integral on a segment
is therefore a positive weighted sum of the two endpoint gaps.
Refining a partition without changing its curve leaves the action
exactly unchanged: the relative spectra on subsegments scale with
length and the integrals add.

\begin{proposition}\label{prop:finitepartition}
On every feasible fixed partition, \eqref{eq:polygonaction} has a unique
minimum among all ordered nodes with the fixed endpoints. If the endpoints
are split over $K$, all nodes of this minimum are split over $K$.
\end{proposition}

\begin{proof}
The objective is nonnegative, lower semicontinuous and convex by
Lemma~\ref{lem:spectralconvexity}, and the potential is affine in
the nodes. Finite cost implies
\[
 d_\infty(c_v(u_j),c_v(u_{j+1}))<\tfrac\pi2\Delta_j.
\]
The triangle inequality from a fixed endpoint thus bounds every node
of every finite sublevel set. A proper lower-semicontinuous convex
function with bounded nonempty sublevel sets on a complete CAT(0)
space attains its infimum. Here one can see attainment directly:
project a fixed origin to the nested closed convex sets with objective
at most the infimum plus $1/m$. The CAT(0) projection inequality gives
\[
 d(p_m,p_k)^2\le d(o,p_k)^2-d(o,p_m)^2\qquad(k\ge m).
\]
The right-hand side tends to zero, so the projections converge to a
point in every sublevel set.

For uniqueness, compare two polygons $c,c'$ on the same partition.
Write $b_{vj}$ for the distance between their $v$-norms at node $j$,
and interpolate their midpoint nodes to obtain $m$.
Decompose $\ell(s/\Delta)$ as
$s^2/(2\Delta^2)$ plus a convex function. Apply
\eqref{eq:fourpoint} to the quadratic part and spectral convexity to
the remainder. The potential at the midpoint nodes is affine, giving
\begin{equation}\label{eq:polygonmidpoint}
 E_I(m)\le\frac{E_I(c)+E_I(c')}{2}
   -\frac1{16}\sum_{v,j}\frac{W_j}{\Delta_j^2}
                         (b_{v,j+1}-b_{vj})^2.
\end{equation}
If both curves minimize, every distance difference vanishes. The
endpoint distances are zero, so every node distance is zero.

If the endpoints are split, apply
Lemma~\ref{lem:finiterealization} to the finite diagram of time nodes,
with identity arrows between consecutive nodes and from $C$ to $A$
at each node. The horizontal spectra preserve the kinetic terms.
The vertical spectra preserve order and the gap $g$, hence the exact
potential. The endpoint norms are retained in full. A split tuple
therefore has the same minimum value, and uniqueness identifies it
with the minimum already constructed.
\end{proof}

Let $\mathcal P_I$ be the finite-action polygons with the fixed endpoints,
and suppose this class is nonempty. Set
$E_I^*=\inf_{c\in\mathcal P_I}E_I(c)$. On continuous ordered curves
with these endpoints define
\begin{equation}\label{eq:localrelaxation}
 \overline E_I(c)=\inf\left\{\liminf_m E_I(c_m):
      c_m\in\mathcal P_I,\ c_m\longrightarrow c\text{ uniformly}\right\},
\end{equation}
with the value $+\infty$ when there is no such sequence.

\begin{proposition}\label{prop:localrelaxedminimum}
The functional \eqref{eq:localrelaxation} has a unique minimum $c_I$,
with value $E_I^*$. Every polygonal minimizing sequence converges
uniformly to $c_I$. Its vertex curves are $\pi/2$-Lipschitz for
$d_\infty$ and $\sqrt n\pi/2$-Lipschitz for $d_2$.
If the endpoints are split, a minimizing sequence can be chosen with
all its nodes split.
\end{proposition}

\begin{proof}
For two finite-action polygons use a common refinement. Since
$W_j\ge\Delta_j$, \eqref{eq:polygonmidpoint}, the zero endpoint
distances and Cauchy--Schwarz give
\begin{equation}\label{eq:localstability}
 \sup_{u\in I,v}d_2(c_v(u),c'_v(u))^2
 \le16|I|\left(\frac{E_I(c)+E_I(c')}{2}-E_I^*\right).
\end{equation}
The node bound extends to every time by Busemann convexity of the
geodesics. Thus any minimizing sequence is uniformly Cauchy in the
complete space of continuous ordered curves. The same inequality
identifies the limits of any two such sequences. The relaxation
attains value $E_I^*$ at this limit, and a diagonal sequence in the
definition of the relaxation proves uniqueness. The barrier bounds
the speed of every finite-action polygon on each segment; the stated
Lipschitz constants pass to the limit. For split endpoints, replace
each polygon by the minimum on its own partition. Its action decreases
and its nodes are split by Proposition~\ref{prop:finitepartition}.
\end{proof}

The final replacement in the proof is an energy comparison. It need
not approximate the original polygon at its interior nodes. It also
does not imply that the uniform limit is split.

\begin{lemma}\label{lem:regularaction}
Suppose $c$ has finitely many fixed-apartment pieces on $I$, with
$C^1$ weights up to their endpoints and all their derivatives in a
compact subinterval of $(-\pi/2,\pi/2)$. Then
\begin{equation}\label{eq:differentialaction}
 \overline E_I(c)=\int_I R(u)
       \left\{\frac12\sum_{v,i}\ell(w'_{vi}(u))+qg(c(u))\right\}\dd u.
\end{equation}
\end{lemma}

\begin{proof}
For an arbitrary norm polygon on a subinterval of length $D$,
\eqref{eq:trianglemajorization} and Jensen's inequality imply
\begin{equation}\label{eq:endpointkineticlower}
 D\sum_i\ell\left(\frac{\sigma_i(h(l),h(r))}{D}\right)
 \le\sum_j\Delta_j\sum_i\ell\left(\frac{\sigma_{ji}}{\Delta_j}\right).
\end{equation}
Consider polygons converging uniformly to $c$. For a fixed subdivision
of $I$, multiply \eqref{eq:endpointkineticlower} on each interval by
its minimum value of $R$ and sum, with the factor $1/2$ and both
vertices. Spectral continuity and lower semicontinuity pass this
endpoint lower bound to the limit. The potential integrals converge
because determinant logarithms are Lipschitz on the compact interval.
Refine the subdivision inside the finitely many apartments of $c$.
Its endpoint spectral quotients are the increments of its weights,
up to permutation, and converge uniformly to their derivatives on
each piece. The strict velocity margin gives convergence of the
lower bounds to the right-hand side of \eqref{eq:differentialaction}.
This proves the lower inequality. This argument also permits converging,
rather than identically fixed, endpoints in the approximating sequence.

Conversely, join the sampled values of $c$ on refining subdivisions
by product geodesics. Within each piece these are the linear
interpolants of the weights. They remain ordered, keep the strict
velocity margin, converge uniformly, and their exact actions converge
to the integral. This proves the upper inequality.
\end{proof}

\subsection{Whole-line comparison curves}
Fix the full boundary triple and set
\[
 r_A(u)=e^{-F(u)}h_A^+,\qquad r_C(u)=e^{F(u)}h_C^+\qquad(u\ge0).
\]
For sufficiently large $u$ this pair is ordered. Let $\calQ$ be the
continuous ordered split curves for which, for some $U>0$,
\begin{equation}\label{eq:exactreferencetails}
 c_A=c_C=h_B\quad(u\le-U),\qquad c_v=r_v\quad(u\ge U),
\end{equation}
and whose restriction to $[-U,U]$ has finitely many fixed-apartment
$C^1$ pieces with strict velocity margins. Derivative jumps at the
finitely many joins are allowed in this comparison class.

This class is nonempty for every triple. Since
$F(U)=(\pi/3)U+O(1)$, the $d_\infty$ distances of $r_A(U),r_C(U)$
from $h_B$ are at most $(\pi/3)U+O(1)$. On $[-U,U]$ the product
geodesic from $(h_B,h_B)$ to $(r_A(U),r_C(U))$ therefore has spectral
speeds less than $\pi/4$ for large $U$. It is ordered by geodesic
convexity. Appending the reference tails gives an element of $\calQ$.

Define the Bregman difference
\[
 D_\ell(v;s)=\ell(v)-\ell(s)-\ell'(s)(v-s),
 \qquad s_A(u)=-\beta(u),\quad s_C(u)=\beta(u).
\]
On $\calQ$ define
\begin{equation}\label{eq:wholelineaction}
 \begin{split}
 J(c)={}&\int_{-\infty}^0 R(u)
   \left\{\frac12\sum_{v,i}\ell(w'_{vi}(u))+qg(c(u))\right\}\dd u\\
 &+\frac12\int_0^\infty R(u)
           \sum_{v,i}D_\ell(w'_{vi}(u);s_v(u))\,\dd u.
 \end{split}
\end{equation}
All integrands are nonnegative, and they vanish outside a compact
interval. The velocity multisets and determinant gap are intrinsic,
so the value does not depend on the apartment presentations.

The renormalization has an exact finite-interval identity. If $J_U$
denotes the two integrals in \eqref{eq:wholelineaction} restricted to
$[-U,U]$, then for every finite regular ordered curve on that interval,
\begin{equation}\label{eq:renormalization}
 J_U(c)=E_{[-U,U]}(c)-\frac{a\sinh U}{2}g(c(U))+C_n(U),
\end{equation}
where the differential action is used and
\[
 C_n(U)=n\int_0^U R(u)
                \{\beta(u)\tan\beta(u)-\ell(\beta(u))\}\,\dd u.
\]
Indeed, $R\ell'(s_C)=a\sinh u=-R\ell'(s_A)$ and
$(a\sinh u)'=aR$. Integration by parts cancels the potential on the
positive half-line. The term at zero vanishes, and the interior
boundary terms cancel by continuity of determinant norms.
No reference endpoint is needed for this truncated identity.
For a curve in $\calQ$ and a cutoff containing its nonreference part,
$J_U(c)=J(c)$. The endpoint term in \eqref{eq:renormalization} then
depends only on the fixed boundary triple and $U$.

\subsection{Existence, uniqueness and the ends}
Let $J_* =\inf_{c\in\calQ}J(c)$. Fix $c_0\in\calQ$ and take the
uniform closure $\calY$ of $\calQ$ among continuous ordered all-norm
curves at finite uniform distance from $c_0$. Use
$\sup_u\max_v d_2$ for this distance. The ambient space, and hence
$\calY$, is complete. Define
\begin{equation}\label{eq:globalrelaxation}
 \overline J(c)=\inf\{\liminf_m J(c_m):c_m\in\calQ,
                                  \ c_m\to c\text{ uniformly}\}.
\end{equation}

\begin{theorem}\label{thm:relaxedminimum}
The functional $\overline J$ has a unique minimum $c_*$ on $\calY$,
with value $J_*$. Every minimizing sequence in $\calQ$ converges
uniformly to $c_*$. Its vertex curves are globally $\pi/2$-Lipschitz
in $d_\infty$. It has the full normalized limits
\eqref{eq:normalizedends} and is equivariant under simultaneous
$\GL_n(K)$ change of basis.
\end{theorem}

\begin{proof}
Take $c,c'\in\calQ$ and a common cutoff $U$. By
Lemma~\ref{lem:regularaction}, they can be approximated on $[-U,U]$
by ordered split polygons on common refinements, with exact endpoints,
uniform convergence and convergence of action. Attach the common
reference tails. For two such polygons retain the coefficients $W_j$
in \eqref{eq:polygonmidpoint}. The midpoint polygon, with reference
tails, belongs to $\calQ$. By \eqref{eq:renormalization}, its action
is at least $J_*$ and the endpoint corrections cancel. Moreover,
\[
 \sum_j\frac{\Delta_j^2}{W_j}
 \le\int_{-U}^U\frac{\dd u}{R(u)}\le\pi.
\]
Weighted Cauchy--Schwarz, starting from the zero endpoint distances,
and Busemann convexity on each segment therefore give, after recovery,
\begin{equation}\label{eq:globalstability}
 \sup_{u,v}d_2(c_v(u),c'_v(u))^2
 \le16\pi\left\{\frac{J(c)+J(c')}{2}-J_*\right\}.
\end{equation}
Thus every minimizing sequence is uniformly Cauchy. Completeness and
the definition of $\calY$ give a limit. The same inequality identifies
all such limits. A diagonal sequence for \eqref{eq:globalrelaxation}
proves that its minimum value is $J_*$ and that the minimum is unique.
The Lipschitz bound follows from the barrier on every regular piece
and passes to the uniform limit.

Since $\ell''\ge1$, one has
$\ell(v)\ge v^2/2$ and $D_\ell(v;s)\ge(v-s)^2/2$.
On a regular apartment piece the metric speed is the Euclidean norm
of the weight derivatives. Integrate this speed from the constant
negative tail, or integrate the speed of the scalar-normalized curve
from the positive reference tail. Weighted Cauchy--Schwarz gives,
for $c\in\calQ$,
\begin{equation}\label{eq:preliminarytailbounds}
 \begin{aligned}
 d_2(c_v(u),h_B)^2
   &\le4J(c)\int_{-\infty}^u\sech s\,\dd s &&(u\le0),\\
 d_2(e^{-\eps_vF(u)}c_v(u),h_v^+)^2
   &\le4J(c)\int_u^\infty\sech s\,\dd s &&(u\ge0),
 \end{aligned}
\end{equation}
where $\eps_A=-1$ and $\eps_C=1$. These estimates pass to a minimizing
sequence with $J(c)$ replaced by $J_*$. They prove the full norm
limits, initially with rate $O(e^{-|u|/2})=O(|t|^{-1/2})$.
Finally, simultaneous change of basis preserves order, relative spectra,
the determinant gap and the reference curves. It preserves $J$ and
its relaxed domain. Uniqueness proves equivariance.
\end{proof}

At this stage $c_*$ is an all-norm curve; splitness and harmonicity
have not been used to obtain its uniqueness. Both will be proved
below, together with the stronger tail rates.

\begin{proposition}\label{prop:finitecutoff}
For large $U$, let $c_U$ be the exact-endpoint relaxed minimum on
$[-U,U]$ with endpoints $(h_B,h_B)$ and $(r_A(U),r_C(U))$, extended
by the reference tails. Put
\[
 j(U)=E_{[-U,U]}^*-\frac{a\sinh U}{2}g(r_A(U),r_C(U))+C_n(U).
\]
Then $j(U)$ decreases to $J_*$, $c_U\in\calY$, and
\begin{equation}\label{eq:finitecutoffconvergence}
 \sup_{u,v}d_2((c_U)_v(u),(c_*)_v(u))^2
       \le8\pi\{j(U)-J_*\}.
\end{equation}
\end{proposition}

\begin{proof}
Let $\calQ_U$ be the curves in $\calQ$ with the exact reference tails
outside $[-U,U]$. The classes increase with $U$ and their union is
$\calQ$. Regular-curve recovery shows that
$j(U)=\inf_{\calQ_U}J$. In the reverse comparison needed for this
equality, an arbitrary all-norm polygon can be replaced by its
split minimum on the same partition, with the same split endpoints
and no larger action. Hence $j(U)\downarrow J_*$.

The split polygonal minimizing sequence for $c_U$, extended by its
reference tails, lies in $\calQ_U$, converges uniformly and has action
tending to $j(U)$. This puts $c_U$ in $\calY$. Apply
\eqref{eq:globalstability} to that sequence and a whole-line
minimizing sequence, then pass to both limits.
\end{proof}

\begin{proposition}\label{prop:extensionminimum}
Canonical scalar extension identifies the whole-line minima over $K$
and over $L$:
\[
 J_K^*=J_L^*,\qquad c_{*,L}(u)=(c_{*,K}(u))_L
 \quad\text{for every }u\in\RR.
\]
\end{proposition}

\begin{proof}
Extension preserves relative spectra, determinant gaps and regular
apartment weights, so it carries $\calQ_K$ into $\calQ_L$ with the
same action. This proves $J_L^*\le J_K^*$.
Conversely, approximate a curve of $\calQ_L$ by polygons on a fixed
cutoff interval, with its exact split $K$ endpoint norms extended to
$L$. The finite formula in the proof of
Lemma~\ref{lem:finiterealization} applies directly to this $L$-diagram,
using the fixed $K$ endpoint matrices and real monomials in $K$ as
parameters. It gives a $K$-polygon preserving all horizontal and
vertical profiles. Its order, exact potential, kinetic action and
full endpoints are therefore preserved. Append the $K$ reference
tails and pass to the infimum. This proves the reverse inequality.

A minimizing sequence over $K$ extends isometrically to a minimizing
sequence over $L$. Its uniform limit is the canonical extension of
$c_{*,K}$. The uniqueness in Theorem~\ref{thm:relaxedminimum}, now
over $L$, identifies that limit with $c_{*,L}$.
\end{proof}
